\documentclass[10pt,a4paper]{article}

\usepackage[spanish,es-nodecimaldot]{babel}
\usepackage[utf8]{inputenc}
\usepackage[T1]{fontenc}
\usepackage{lmodern}
\usepackage[margin=1.35cm]{geometry}
\usepackage[hidelinks]{hyperref}
\usepackage{enumitem}
\usepackage{titlesec}
\usepackage{tabularx}
\usepackage{xcolor}
\usepackage{array}

\definecolor{heading}{HTML}{111827}
\definecolor{muted}{HTML}{4B5563}
\definecolor{rulecolor}{HTML}{D1D5DB}

\pagestyle{empty}
\setlength{\parindent}{0pt}
\setlength{\parskip}{2pt}
\hyphenpenalty=10000
\exhyphenpenalty=10000
\emergencystretch=2em
\setlist[itemize]{leftmargin=1.2em, noitemsep, topsep=0pt, parsep=0pt, partopsep=0pt}

\titleformat{\section}
  {\large\bfseries\color{heading}}
  {}
  {0pt}
  {\MakeUppercase}
  [\vspace{-2pt}\color{rulecolor}\titlerule]
\titlespacing*{\section}{0pt}{8pt}{8pt}

\newcommand{\cvheader}[6]{
  \begin{center}
    {\LARGE\bfseries\MakeUppercase{#1}}\\[3pt]
    {\color{muted}#2}\\[3pt]
    {#3 \textbullet{} \href{mailto:#4}{#4} \textbullet{} \href{https://#5}{#5} \textbullet{} \href{https://#6}{#6}}
  \end{center}
  \vspace{-4pt}
}

\newcommand{\experienceentry}[5]{
  \begingroup
  \setlength{\parskip}{0pt}
  \textbf{#1} \hfill {\color{muted}#3}\par
  \textbf{#2} \hfill {\color{muted}#4}
  #5
  \endgroup
  \vspace{4pt}
}

\newcommand{\projectentry}[5]{
  \textbf{#1} \hfill {\color{muted}#2}\\
  #3\\[-1pt]
  \textbf{Tecnologías:} #4\\[-1pt]
  \textbf{Repositorio:} \href{https://#5}{#5}\vspace{2pt}
}

\begin{document}

\cvheader
  {Facundo Joel Rodriguez}
  {Licenciado en Sistemas de Información}
  {Misiones, Argentina}
  {facurodri00@gmail.com}
  {linkedin.com/in/facundo-joel-rodriguez}
  {github.com/facurodrij}

\section{Perfil}
Desarrollador Full Stack y Licenciado en Sistemas de Información. Construyo backends y aplicaciones web con Python y Node.js, y me ocupo de que lleguen a producción: APIs, bases de datos, CI/CD y despliegue en la nube. Actualmente desarrollo en Sherpa la plataforma que automatiza la creación de chatbots de WhatsApp para productores y brokers de seguros. Trabajo con metodologías ágiles, priorizando la calidad del código, la mantenibilidad y el trabajo colaborativo.

\section{Experiencia}
\experienceentry
  {Sherpa.wtf}
  {Desarrollador Full Stack}
  {Remoto}
  {Marzo 2025 -- Presente}
  {\begin{itemize}
    \item Desarrollo y mantengo Genesis, plataforma que automatiza la creación de chatbots para brokers de seguros.
    \item Participé en el despliegue y la migración de la plataforma hacia el CRM propio de Sherpa.
    \item Contribuí a la aprobación de la app en Meta para la API de WhatsApp Business y el Registro Insertado.
    \item Trabajo sobre infraestructura en AWS definida con Terraform y servicios en contenedores Docker.
  \end{itemize}
  \textbf{Tecnologías:} Node.js, TypeScript, WhatsApp Business API, AWS, Docker, Terraform}

\experienceentry
  {Repuestos Stock Car}
  {Desarrollador Full Stack | Profesional independiente}
  {Posadas, Misiones}
  {Julio 2024 -- Febrero 2025}
  {\begin{itemize}
    \item Diseñé y construí de punta a punta el sistema de facturación e inventario de la empresa, que hoy sigue en producción.
    \item Relevé los requerimientos con el cliente y definí la arquitectura del sistema.
    \item Brindé soporte técnico y mejoras continuas para asegurar continuidad operativa y satisfacción del cliente.
  \end{itemize}
  \textbf{Tecnologías:} Python, Flask, React}

\experienceentry
  {Argentina Programa 4.0}
  {Tutor Docente}
  {Remoto}
  {Junio 2023 -- Septiembre 2023}
  {\begin{itemize}
    \item Acompañé estudiantes en IA con Python en el trayecto Programador Junior en Machine Learning.
  \end{itemize}}

\section{Educación}
\textbf{Universidad Nacional de Misiones} \hfill {\color{muted}Apóstoles, Misiones}\\
\textbf{Licenciado en Sistemas de Información} \hfill {\color{muted}Feb. 2019 -- Jun. 2024}\\
Promedio sin aplazos: 7.93/10\\
Trabajo Final: Tesis de Grado con Tecnología Blockchain\\[2pt]
\textbf{Analista en Sistemas de Computación} \hfill {\color{muted}Feb. 2019 -- Mar. 2023}\\
Promedio sin aplazos: 7.77/10\\
Proyecto Final: Sistema de Gestión para un Club Deportivo

\section{Proyectos destacados}
\projectentry
  {Tesis de Grado con Tecnología Blockchain \textnormal{(Académico)}}
  {2024}
  {Prototipo basado en Blockchain para el Tribunal de Cuentas de la Provincia de Misiones, con el objetivo de garantizar la integridad y confiabilidad de las presentaciones digitales de rendiciones de cuentas municipales.}
  {Python, Django, jQuery, Bootstrap, SQLite, Solidity, Truffle, Ganache}
  {github.com/facurodrij/TrabajoFinal\_LSI}

\projectentry
  {Sistema de Gestión para un Club Deportivo \textnormal{(Académico)}}
  {2023}
  {Software de gestión para un club deportivo que abarca la administración de socios, eventos, ventas de tickets y reservas de canchas. La integración de la API de Mercado Pago permitió simular la automatización de pagos en línea, optimizando procesos administrativos y mejorando la experiencia del usuario en pruebas funcionales.}
  {Python, Django, PostgreSQL, jQuery, Bootstrap}
  {github.com/facurodrij/TrabajoFinal\_ASC}

\section{Habilidades}
\textbf{Técnicas:} Python, Node.js, TypeScript, Django, Flask, Express, React, Next.js, REST APIs\\
\textbf{Bases de Datos:} SQLite, PostgreSQL, MongoDB\\
\textbf{Cloud \& DevOps:} Docker, AWS, GCP, Terraform, GitHub Actions, Linux\\
\textbf{Herramientas:} Git/GitHub, Postman, Cursor, Jira, Notion\\
\textbf{Metodologías:} Scrum, Prácticas Ágiles, UP, XP, Desarrollo iterativo e incremental\\
\textbf{Idiomas:} Español nativo, Inglés B1 \textit{intermediate}

\end{document}
